\documentclass[a4paper,12pt]{article}
\usepackage[utf8]{inputenc}
\usepackage[T1]{fontenc}
\usepackage[ngerman]{babel}
\usepackage{amsmath, amssymb}
\usepackage{geometry}
\usepackage{parskip}
\usepackage{titlesec}
\usepackage{enumitem}
\usepackage{array}

\geometry{left=3cm, right=3cm, top=2.5cm, bottom=2.5cm}

\titleformat{\section}{\large\bfseries}{}{0em}{}[\titlerule]
\titleformat{\subsection}{\normalsize\bfseries\scshape}{}{0em}{}

% Glossar-Eintrag: Nummer, fetter Titel, Inhalt
\newcommand{\gentry}[3]{%
  \noindent\textbf{#1\quad #2}\nopagebreak\\[0.2em]
  #3\par\smallskip
}

\begin{document}

\begin{center}
  {\LARGE\bfseries Glossar zu Lektion 3}\\[0.5em]
  {\large Mathematische Grundlagen (Modul 61111)}
\end{center}

\vspace{1em}

Dieses Glossar fasst die wesentlichen Definitionen, Merkregeln, Sätze und
Propositionen der dritten Lektion kompakt zusammen. Nummerierungen
entsprechen dem Lehrtext.

% ======================================================================
\section*{7.1 \quad Lineare Unabhängigkeit}

\gentry{7.1.1}{Definition (Triviale / nicht-triviale Darstellung)}{%
  Sei $V$ ein $K$-Vektorraum und $v_1,\ldots,v_n \in V$.
  \begin{itemize}[topsep=2pt,itemsep=1pt]
    \item Die Linearkombination $\sum_{i=1}^n a_i v_i = 0$ heißt
          \textbf{triviale Darstellung} des Nullvektors, wenn $a_i = 0$ für alle $i$.
    \item Sie heißt \textbf{nicht-triviale Darstellung}, wenn mindestens ein
          $a_i \neq 0$ ist.
  \end{itemize}
}

\gentry{7.1.2}{Definition (Linear unabhängig / linear abhängig)}{%
  Vektoren $v_1,\ldots,v_n \in V$ heißen \textbf{linear unabhängig}, falls
  \[
    \sum_{i=1}^n a_i v_i = 0 \;\Longrightarrow\; a_i = 0 \text{ für alle } i.
  \]
  Sie heißen \textbf{linear abhängig}, falls eine nicht-triviale Darstellung
  des Nullvektors durch $v_1,\ldots,v_n$ existiert.
  \smallskip\\
  \textbf{Spezialfälle:} Ein einzelner Vektor $v$ ist linear unabhängig genau
  dann, wenn $v \neq 0$. Zwei Vektoren $v_1, v_2$ sind linear abhängig genau
  dann, wenn einer ein skalares Vielfaches des anderen ist.
}

\gentry{7.1.6}{Proposition (Charakterisierung linear abhängiger Vektoren)}{%
  $v_1,\ldots,v_n \in V$ sind linear abhängig genau dann, wenn es ein $i$
  gibt, so dass
  \[
    v_i \in \langle v_1,\ldots,v_{i-1},v_{i+1},\ldots,v_n \rangle.
  \]
  \textbf{Kriterium:} $v_1,\ldots,v_n$ sind linear unabhängig genau dann,
  wenn die Koeffizientenmatrix $\begin{pmatrix}v_1 & \cdots & v_n\end{pmatrix}$
  (Spalten = Vektoren) den Rang $n$ hat.
}

% ======================================================================
\section*{7.2 \quad Basen endlich erzeugter Vektorräume}

\gentry{7.2.1}{Definition (Basis)}{%
  Sei $V$ ein endlich erzeugter $K$-Vektorraum, $V \neq \{0\}$. Die Vektoren
  $v_1,\ldots,v_n \in V$ heißen eine \textbf{Basis} von $V$, falls sie ein
  Erzeugendensystem von $V$ sind und linear unabhängig sind.
  \smallskip\\
  Die \textbf{leere Menge} $\emptyset$ ist per Definition die Basis von $\{0\}$.
}

\gentry{7.2.4--7.2.6}{Standardbasen}{%
  \begin{itemize}[topsep=2pt,itemsep=1pt]
    \item $K^n$: Einheitsvektoren $e_1,\ldots,e_n$ (kanonische Basis).
    \item $M_{mn}(K)$: Matrizen $E_{ij}$, $1 \le i \le m$, $1 \le j \le n$.
    \item Polynome vom Grad $\le n$: $e_0 = 1,\, e_1 = T,\ldots,\, e_n = T^n$.
  \end{itemize}
}

\gentry{7.2.7}{Proposition (Charakterisierung von Basen)}{%
  $v_1,\ldots,v_n$ ist eine Basis von $V$ genau dann, wenn sich jedes
  $v \in V$ \textbf{eindeutig} als Linearkombination $v = \sum_{i=1}^n a_i v_i$
  schreiben lässt.
}

\gentry{7.2.8}{Proposition (Existenz von Basen)}{%
  Sei $V \neq \{0\}$ endlich erzeugt.
  \begin{enumerate}[label=(\alph*),topsep=2pt,itemsep=1pt]
    \item Ist $v_1,\ldots,v_n$ ein Erzeugendensystem und gilt
          $v_i \in \langle v_1,\ldots,\hat{v}_i,\ldots,v_n\rangle$, so ist
          $v_1,\ldots,\hat{v}_i,\ldots,v_n$ ebenfalls ein Erzeugendensystem.
    \item Jedes endliche Erzeugendensystem enthält eine Basis.
  \end{enumerate}
}

\gentry{7.2.11}{Bemerkung}{%
  Sei $V$ ein $K$-Vektorraum und $v_1,\ldots,v_n \in V$.
  \begin{enumerate}[label=(\alph*),topsep=2pt,itemsep=1pt]
    \item Sind $v_1,\ldots,v_n$ linear unabhängig und gilt
          $v_{n+1} \notin \langle v_1,\ldots,v_n\rangle$, so sind
          $v_1,\ldots,v_{n+1}$ linear unabhängig.
    \item Ist $v_1,\ldots,v_n$ ein Erzeugendensystem von $V$, so sind
          $v_1,\ldots,v_n,v_{n+1}$ für alle $v_{n+1} \in V$ linear abhängig.
  \end{enumerate}
}

% ======================================================================
\section*{7.3 \quad Der Austauschsatz von Steinitz}

\gentry{7.3.1}{Lemma (Austauschlemma)}{%
  Sei $u_1,\ldots,u_n$ eine Basis von $V$ und sei $v \in V$, $v \neq 0$. Dann
  gibt es ein $i$ mit $1 \le i \le n$, so dass
  \[
    u_1,\ldots,u_{i-1},\; v,\; u_{i+1},\ldots,u_n
  \]
  eine Basis von $V$ ist. (Austausch: Der Basisvektor $u_i$ wird durch $v$
  ersetzt, wenn der Koeffizient von $u_i$ in der Darstellung von $v$ bzgl.\ der
  Basis ungleich Null ist.)
}

\gentry{7.3.3}{Satz (Austauschsatz von Steinitz)}{%
  Sei $u_1,\ldots,u_n$ eine Basis von $V$ und seien $v_1,\ldots,v_m$ linear
  unabhängige Vektoren in $V$. Dann gilt $n \ge m$, und es gibt Vektoren
  $u_{i_{m+1}},\ldots,u_{i_n} \in \{u_1,\ldots,u_n\}$, so dass
  \[
    v_1,\ldots,v_m,\; u_{i_{m+1}},\ldots,u_{i_n}
  \]
  eine Basis von $V$ ist.
}

\gentry{7.3.4}{Korollar (Charakterisierung endlich erzeugter Vektorräume)}{%
  $V$ ist endlich erzeugt genau dann, wenn es ein $n \in \mathbb{N}$ gibt, so
  dass jede Menge von linear unabhängigen Vektoren in $V$ höchstens $n$
  Elemente besitzt.
}

\gentry{7.3.5}{Korollar (Basisergänzungssatz)}{%
  Sei $V$ endlich erzeugt und seien $w_1,\ldots,w_m$ linear unabhängig in $V$.
  Dann lassen sich $w_1,\ldots,w_m$ zu einer Basis von $V$ ergänzen: Es gibt
  $v_{m+1},\ldots,v_n$, so dass $w_1,\ldots,w_m,v_{m+1},\ldots,v_n$ eine Basis
  von $V$ ist.
}

\gentry{7.3.6}{Korollar (Eindeutigkeit der Basismächtigkeit)}{%
  Je zwei Basen eines endlich erzeugten Vektorraums $V$ haben die gleiche
  Anzahl von Vektoren.
}

% ======================================================================
\section*{7.4 \quad Dimension}

\gentry{7.4.1}{Definition (Dimension)}{%
  Sei $V$ ein $K$-Vektorraum.
  \begin{enumerate}[label=(\alph*),topsep=2pt,itemsep=1pt]
    \item $\dim_K(V) = \infty$, falls $V$ nicht endlich erzeugt ist.
    \item $\dim_K(V) = n$, falls $V \neq \{0\}$ endlich erzeugt ist und jede
          Basis aus $n$ Vektoren besteht.
    \item $\dim_K(\{0\}) = 0$.
  \end{enumerate}
  \textbf{Beispiele:} $\dim(K^n) = n$,\; $\dim(M_{mn}(K)) = mn$,\;
  $\dim(\text{Pol}_{\le n}) = n+1$.
}

\gentry{7.4.4}{Proposition}{%
  Sei $\dim(V) = n < \infty$.
  \begin{enumerate}[label=(\alph*),topsep=2pt,itemsep=1pt]
    \item $n$ linear unabhängige Vektoren in $V$ bilden bereits eine Basis.
    \item Ein Erzeugendensystem von $V$ mit genau $n$ Vektoren ist bereits eine
          Basis.
  \end{enumerate}
  \textbf{Merke:} Bei bekannter Dimension genügt der Nachweis \emph{einer}
  der beiden Basiseigenschaften.
}

\gentry{7.4.5}{Korollar (Dimensionsformel für Unterräume)}{%
  Sei $U$ ein Unterraum des endlich erzeugten Vektorraums $V$. Dann gilt
  \[
    \dim(U) \le \dim(V),
  \]
  und Gleichheit gilt genau dann, wenn $U = V$.
}

\gentry{7.4.6}{Proposition (Dimensionsformel für Summe und Durchschnitt)}{%
  Seien $U$ und $W$ endlichdimensionale Unterräume von $V$. Dann gilt
  \[
    \dim(U + W) = \dim(U) + \dim(W) - \dim(U \cap W).
  \]
  \textbf{Beweisskizze:} Basis $x_1,\ldots,x_r$ von $U \cap W$ durch
  $u_1,\ldots,u_s$ zu Basis von $U$ und durch $w_1,\ldots,w_t$ zu Basis von
  $W$ ergänzen; dann ist $x_1,\ldots,x_r,u_1,\ldots,u_s,w_1,\ldots,w_t$ eine
  Basis von $U+W$.
}

% ======================================================================
\section*{8.1 \quad Lineare Abbildungen: Definition und Eigenschaften}

\gentry{8.1.1}{Definition (Lineare Abbildung / Vektorraumhomomorphismus)}{%
  Seien $V, W$ $K$-Vektorräume. Eine Abbildung $f: V \to W$ heißt
  \textbf{(K-)linear}, falls für alle $v_1,v_2 \in V$ und $a \in K$ gilt:
  \[
    f(v_1 + v_2) = f(v_1) + f(v_2), \qquad f(av) = af(v).
  \]
  Äquivalent: $f(\sum a_i v_i) = \sum a_i f(v_i)$ (Bild~8.1.11).
}

\gentry{8.1.4}{Bemerkung (Elementare Eigenschaften)}{%
  Sei $f: V \to W$ linear. Dann gilt:
  \begin{enumerate}[label=(\alph*),topsep=2pt,itemsep=1pt]
    \item $f(0_V) = 0_W$.
    \item $f(-v) = -f(v)$ für alle $v \in V$.
  \end{enumerate}
}

\gentry{8.1.5}{Definition (Isomorphismus)}{%
  Eine bijektive lineare Abbildung $f: V \to W$ heißt \textbf{Isomorphismus}.
  Man schreibt $V \cong W$ und nennt $V$ und $W$ \textbf{isomorph}.
}

\gentry{8.1.8}{Proposition}{%
  Ist $f: V \to W$ ein Isomorphismus, so ist auch $f^{-1}: W \to V$ ein
  Isomorphismus.
}

\gentry{8.1.9}{Proposition (Komposition)}{%
  Seien $f: U \to V$ und $g: V \to W$ lineare Abbildungen. Dann ist
  $g \circ f: U \to W$ linear. Sind $f$ und $g$ Isomorphismen, so auch
  $g \circ f$ (Kor.~8.1.10).
}

% ======================================================================
\section*{8.2 \quad Isomorphe Vektorräume}

\gentry{8.2.2}{Satz (Struktur endlich erzeugter Vektorräume)}{%
  Sei $V$ endlich erzeugt und sei $B = (v_1,\ldots,v_n)$ eine Basis von $V$.
  Die Abbildung
  \[
    \kappa_B: V \to K^n, \quad v = \textstyle\sum_{i=1}^n a_i v_i
    \;\mapsto\; \begin{pmatrix}a_1\\\vdots\\a_n\end{pmatrix}
  \]
  ist ein Isomorphismus.
}

\gentry{8.2.4}{Definition (Koordinatenvektor)}{%
  Sei $B$ eine Basis von $V$ und $v \in V$. Dann heißt $\kappa_B(v) \in K^n$
  der \textbf{Koordinatenvektor} von $v$ bezüglich $B$.
}

\gentry{8.2.6}{Korollar}{%
  Seien $V$ und $W$ endlich erzeugte $K$-Vektorräume.
  \begin{enumerate}[label=(\alph*),topsep=2pt,itemsep=1pt]
    \item $V \cong W \;\Leftrightarrow\; \dim(V) = \dim(W)$.
    \item $\dim(V) = n \;\Rightarrow\; V \cong K^n$.
  \end{enumerate}
}

% ======================================================================
\section*{8.3 \quad Kern und Bild}

\gentry{8.3.1}{Definition (Bild)}{%
  Sei $f: V \to W$ linear. Das \textbf{Bild} von $f$ ist
  \[
    \mathrm{Bild}(f) = \{w \in W \mid \exists\, v \in V: f(v) = w\}.
  \]
  $\mathrm{Bild}(f)$ ist ein Unterraum von $W$ (Bem.~8.3.2).
  Ist $v_1,\ldots,v_n$ eine Basis von $V$, so ist $f(v_1),\ldots,f(v_n)$
  ein Erzeugendensystem von $\mathrm{Bild}(f)$ (Prop.~8.3.4).
}

\gentry{8.3.5}{Korollar}{%
  Sei $v_1,\ldots,v_n$ eine Basis von $V$. Dann ist $f$ surjektiv genau dann,
  wenn $f(v_1),\ldots,f(v_n)$ ein Erzeugendensystem von $W$ ist.
}

\gentry{8.3.7}{Definition (Kern)}{%
  Sei $f: V \to W$ linear. Der \textbf{Kern} von $f$ ist
  \[
    \mathrm{Kern}(f) = \{v \in V \mid f(v) = 0\}.
  \]
  $\mathrm{Kern}(f)$ ist ein Unterraum von $V$ (Bem.~8.3.8).
}

\gentry{8.3.10}{Proposition}{%
  $f: V \to W$ ist injektiv genau dann, wenn $\mathrm{Kern}(f) = \{0\}$.
}

\gentry{8.3.12}{Definition (Rang einer linearen Abbildung)}{%
  $\mathrm{Rg}(f) := \dim(\mathrm{Bild}(f))$.
}

\gentry{8.3.13}{Satz}{%
  Sei $\dim(V) = n$, sei $v_1,\ldots,v_r$ eine Basis von $\mathrm{Kern}(f)$ und
  sei $v_1,\ldots,v_r,v_{r+1},\ldots,v_n$ eine Basis von $V$. Dann ist
  $f(v_{r+1}),\ldots,f(v_n)$ eine Basis von $\mathrm{Bild}(f)$.
}

\gentry{8.3.14}{Korollar (Rangsatz)}{%
  Sei $V$ endlich erzeugt und $f: V \to W$ linear. Dann gilt
  \[
    \dim(V) = \dim(\mathrm{Kern}(f)) + \mathrm{Rg}(f).
  \]
}

\gentry{8.3.17}{Korollar}{%
  Seien $V, W$ endlich erzeugt mit $\dim(V) = \dim(W)$, und sei $f: V \to W$
  linear. Dann gilt:
  \[
    f \text{ surjektiv} \;\Longleftrightarrow\; f \text{ injektiv}
    \;\Longleftrightarrow\; f \text{ bijektiv (Isomorphismus)}.
  \]
}

\gentry{8.3.18}{Korollar}{%
  Sei $v_1,\ldots,v_n$ eine Basis von $V$ und $f: V \to W$ linear.
  \begin{enumerate}[label=(\alph*),topsep=2pt,itemsep=1pt]
    \item $f$ injektiv $\Leftrightarrow$ $f(v_1),\ldots,f(v_n)$ Basis von $\mathrm{Bild}(f)$.
    \item $f$ surjektiv $\Leftrightarrow$ $f(v_1),\ldots,f(v_n)$ Erzeugendensystem von $W$.
    \item $f$ bijektiv $\Leftrightarrow$ $f(v_1),\ldots,f(v_n)$ Basis von $W$.
  \end{enumerate}
}

\gentry{8.3.19}{Satz (Struktursatz endlich erzeugter Vektorräume)}{%
  Seien $V, W$ endlich erzeugte $K$-Vektorräume.
  \begin{enumerate}[label=(\alph*),topsep=2pt,itemsep=1pt]
    \item $V \cong W \;\Leftrightarrow\; \dim(V) = \dim(W)$.
    \item $\dim(V) = n \;\Rightarrow\; V \cong K^n$.
  \end{enumerate}
}

% ======================================================================
\section*{8.4 \quad Lineare Abbildungen und Basen}

\gentry{8.4.1}{Satz (Beschreibung linearer Abbildungen durch Bilder einer Basis)}{%
  Sei $v_1,\ldots,v_n$ eine Basis von $V$ und seien $w_1,\ldots,w_n \in W$
  beliebig. Dann gibt es \textbf{genau eine} lineare Abbildung $f: V \to W$
  mit $f(v_i) = w_i$ für alle $1 \le i \le n$.
  \smallskip\\
  \textbf{Merke:} Eine lineare Abbildung ist vollständig durch ihre Werte auf
  einer Basis bestimmt.
}

\gentry{8.4.6}{Korollar}{%
  Sei $\dim(V) = \dim(W) = n$, und seien $v_1,\ldots,v_n$ bzw.\ $w_1,\ldots,w_n$
  Basen von $V$ bzw.\ $W$. Dann gibt es genau einen Isomorphismus $f: V \to W$
  mit $f(v_i) = w_i$.
}

\gentry{8.4.8}{Definition (Basiswechsel)}{%
  Sei $V$ endlich erzeugt und seien $B, B'$ zwei Basen von $V$. Der eindeutige
  Isomorphismus, der $B$ in $B'$ überführt, heißt \textbf{Basiswechsel} (oder
  Basistransformation) von $B$ nach $B'$.
}

% ======================================================================
\section*{9.1 \quad Der Vektorraum $\mathrm{Hom}_K(V,W)$}

\gentry{9.0.1}{Definition (System von Vektoren)}{%
  Ein \textbf{System} von $n$ Vektoren $(v_1,\ldots,v_n)$ in $V$ ist ein
  geordnetes $n$-Tupel. Im Gegensatz zur Menge ist die Reihenfolge wesentlich.
  Basen werden als Systeme notiert.
}

\gentry{9.1.1}{Definition ($\mathrm{Hom}_K(V,W)$)}{%
  $\mathrm{Hom}_K(V,W)$ ist die Menge aller linearen Abbildungen von $V$ nach
  $W$, versehen mit punktweiser Addition $(f+g)(v) = f(v)+g(v)$ und
  Skalarmultiplikation $(af)(v) = af(v)$.
}

\gentry{9.1.2}{Proposition}{%
  $\mathrm{Hom}_K(V,W)$ ist ein $K$-Vektorraum.
  Das neutrale Element ist die Nullabbildung $0: v \mapsto 0$.
}

\gentry{9.1.4}{Definition (Matrixdarstellung)}{%
  Seien $B = (v_1,\ldots,v_n)$ eine Basis von $V$ und
  $C = (w_1,\ldots,w_m)$ eine Basis von $W$, und sei $f: V \to W$ linear.
  Schreibe $f(v_j) = \sum_{i=1}^m a_{ij} w_i$. Die \textbf{Matrixdarstellung}
  von $f$ bezüglich $B$ und $C$ ist
  \[
    {}_C M_B(f) = (a_{ij}) \in M_{mn}(K),
  \]
  wobei die $j$-te Spalte der Koordinatenvektor von $f(v_j)$ bzgl.\ $C$ ist.
}

\gentry{9.1.9}{Satz (Darstellung linearer Abbildungen durch Matrizen)}{%
  Die Abbildung ${}_C M_B: \mathrm{Hom}_K(V,W) \to M_{mn}(K)$ ist ein
  Isomorphismus. Insbesondere ist jeder $m \times n$-Matrix $A$ genau eine
  lineare Abbildung $f_A$ zugeordnet.
}

\gentry{9.1.10}{Korollar (Dimensionsformel für Homomorphismenräume)}{%
  \[
    \dim(\mathrm{Hom}_K(V,W)) = \dim(V) \cdot \dim(W).
  \]
}

% ======================================================================
\section*{9.2 \quad Koordinatenvektoren und Matrixdarstellungen}

\gentry{9.2.1}{Proposition}{%
  Seien $B$ Basis von $V$, $C$ Basis von $W$ und $f: V \to W$ linear.
  Dann gilt für alle $v \in V$:
  \[
    {}_C M_B(f) \cdot \kappa_B(v) = \kappa_C(f(v)).
  \]
  \textbf{Merke:} Multiplikation der Matrixdarstellung mit dem Koordinatenvektor
  liefert den Koordinatenvektor des Bildvektors.
}

\gentry{9.2.3}{Korollar}{%
  $\mathrm{Kern}(f)$ ist isomorph zur Lösungsmenge $U$ des homogenen linearen
  Gleichungssystems ${}_C M_B(f)\,x = 0$. Insbesondere:
  $\dim(\mathrm{Kern}(f)) = \dim(U) = n - \mathrm{Rg}({}_C M_B(f))$.
}

\gentry{9.2.4}{Korollar}{%
  $\mathrm{Rg}(f) = \mathrm{Rg}({}_C M_B(f))$.
  Rang der linearen Abbildung = Rang ihrer Matrixdarstellung.
}

\gentry{9.2.5}{Korollar}{%
  Sei $A \in M_{mn}(K)$ mit Spalten $v_1,\ldots,v_n$. Dann gilt
  \[
    \mathrm{Rg}(A) = \dim\langle v_1,\ldots,v_n \rangle.
  \]
  Zeilenrang und Spaltenrang einer Matrix stimmen überein.
}

% ======================================================================
\section*{9.3 \quad Matrizenprodukt und Komposition}

\gentry{9.3.1}{Proposition (Matrizenmultiplikation = Komposition)}{%
  Seien $f: V \to W$ und $g: W \to X$ linear, und seien $B, C, D$ Basen von
  $V, W$ bzw.\ $X$. Dann gilt
  \[
    {}_D M_B(g \circ f) = {}_D M_C(g) \cdot {}_C M_B(f).
  \]
  \textbf{Merke:} Matrizenmultiplikation entspricht der Komposition linearer
  Abbildungen.
}

\gentry{9.3.2}{Aufgabe/Korollar (Inverse und Matrixdarstellung)}{%
  Sei $f: V \to W$ ein Isomorphismus. Dann gilt
  \[
    {}_C M_B(f)^{-1} = {}_B M_C(f^{-1}).
  \]
  Invertierbarkeit der Matrixdarstellung entspricht dem Isomorphismus.
}

\gentry{9.3.3}{Korollar}{%
  Seien $A' \in M_{ms}(K)$, $A'' \in M_{sn}(K)$ und $A = A'A''$. Dann gilt
  \[
    \mathrm{Rg}(A) \le \mathrm{Rg}(A') \quad\text{und}\quad
    \mathrm{Rg}(A) \le \mathrm{Rg}(A'').
  \]
}

\end{document}
